\documentclass{article}

\usepackage[utf8]{inputenc}
\usepackage[T1]{fontenc}
\usepackage{graphicx}
\usepackage{float}
\usepackage{amsmath}
\usepackage{amssymb}
\usepackage{hyperref}

\title{\textbf{Equação de Advecção-Difusão}\\
Modelação em Física e Engenharia}

\author{Tiago Bernardo - 53117 \and Gabriel Marques - 53087}

\date{Maio 2026}

\begin{document}

\maketitle

\section{Introdução}

A equação de advecção--difusão constitui uma das equações diferenciais parciais fundamentais na modelação de fenómenos de transporte em Física e Engenharia.

Esta equação descreve a evolução temporal e espacial de uma grandeza sujeita simultaneamente a transporte por advecção e a processos de difusão.

O objetivo deste trabalho consiste na resolução numérica da equação de advecção--difusão unidimensional através do método das diferenças finitas, estudando a evolução de um perfil inicial Gaussiano em diferentes regimes físicos caracterizados pelo número de Péclet.

Adicionalmente, pretende-se comparar os esquemas numéricos Upwind e Lax--Wendroff, analisando estabilidade, precisão e difusão numérica.

\section{Fundamentos Teóricos}

A equação de advecção--difusão unidimensional é dada por:

\[
\frac{\partial u}{\partial t}
+
v \frac{\partial u}{\partial x}
=
D \frac{\partial^2 u}{\partial x^2}.
\]

O termo advectivo está associado ao transporte da solução, enquanto o termo difusivo representa o seu espalhamento.

Os parâmetros adimensionais utilizados são:

\[
C = \frac{v \Delta t}{\Delta x},
\]

\[
r = \frac{D \Delta t}{(\Delta x)^2}.
\]

O número de Péclet define a razão entre advecção e difusão:

\[
Pe = \frac{v L}{D}.
\]

Valores elevados de \(Pe\) indicam regime dominado pela advecção, enquanto valores reduzidos indicam domínio da difusão.

Como condição inicial foi utilizado um perfil Gaussiano:

\[
u(x,0) = \exp\left(-\frac{(x - x_0)^2}{\sigma^2}\right).
\]

\section{Métodos Numéricos}

A discretização temporal foi feita pelo método de Euler explícito:

\[
\frac{\partial u}{\partial t}
\approx
\frac{u_i^{n+1} - u_i^n}{\Delta t}.
\]

O termo difusivo foi discretizado por diferenças centradas:

\[
\frac{\partial^2 u}{\partial x^2}
\approx
\frac{u_{i+1}^n - 2u_i^n + u_{i-1}^n}{(\Delta x)^2}.
\]

Foram usadas condições de fronteira de Dirichlet homogéneas:

\[
u(0,t) = 0, \quad u(L,t) = 0.
\]

\subsection{Esquema Upwind}

O esquema Upwind utiliza:

\[
\frac{\partial u}{\partial x}
\approx
\frac{u_i^n - u_{i-1}^n}{\Delta x}.
\]

Este esquema é robusto e estável, mas introduz difusão numérica artificial, reduzindo a precisão em regimes dominados pela advecção.

\subsection{Esquema de Lax--Wendroff}

O esquema de Lax--Wendroff é dado por:

\[
u_i^{n+1}
=
u_i^n
-
\frac{C}{2} \left(u_{i+1}^n - u_{i-1}^n\right)
+
\frac{C^2}{2} \left(u_{i+1}^n - 2u_i^n + u_{i-1}^n\right)
+
r \left(u_{i+1}^n - 2u_i^n + u_{i-1}^n\right).
\]

Este método apresenta maior precisão e menor difusão numérica.

\subsection{Estabilidade Numérica}

A estabilidade da solução requer:

\[
C \le 1,
\qquad
r \le \frac{1}{2}.
\]

Todos os casos estudados respeitam estas condições.

\section{Resultados Numéricos}

Foram estudados três regimes físicos distintos:

\begin{itemize}
\item Advecção dominante;
\item Difusão dominante;
\item Regime equilibrado.
\end{itemize}

\begin{figure}[H]
\centering
\includegraphics[width=0.4\textwidth]{parametros.png}
\caption{Parâmetros adimensionais utilizados nas simulações.}
\end{figure}

\subsection{Número de Péclet}

\begin{table}[H]
\centering
\begin{tabular}{cccc}
\hline
Regime & $v$ & $D$ & $Pe$ \\
\hline
Advecção dominante & 2.0 & $10^{-5}$ & 200000 \\
Equilibrado & 1.0 & $5 \times 10^{-3}$ & 200 \\
Difusão dominante & 0.05 & $10^{-2}$ & 5 \\
\hline
\end{tabular}
\caption{Número de Péclet nos diferentes regimes estudados.}
\end{table}

\subsection{Comparação entre esquemas numéricos}

\begin{figure}[H]
\centering
\begin{minipage}{0.49\textwidth}
    \centering
    \includegraphics[width=\textwidth]{Gráficos/advecção_dominante_upwind.png}
\end{minipage}
\hfill
\begin{minipage}{0.49\textwidth}
    \centering
    \includegraphics[width=\textwidth]{Gráficos/advecção_dominante.png}
\end{minipage}

\caption{Comparação entre os esquemas Upwind (esquerda) e Lax--Wendroff (direita) no regime advectivo.}
\end{figure}

O esquema Upwind apresenta maior difusão numérica, conduzindo a perda de amplitude e suavização do perfil. O esquema de Lax--Wendroff preserva melhor a forma da solução e apresenta maior precisão.

\subsection{Regime dominado pela advecção}

Neste regime observa-se transporte dominante do perfil Gaussiano.

\begin{figure}[H]
\centering
\includegraphics[width=0.8\textwidth]{Gráficos/advecção_dominante.png}
\caption{Regime dominado pela advecção.}
\end{figure}

\subsection{Regime dominado pela difusão}

Neste caso verifica-se alargamento progressivo do perfil e redução da amplitude máxima.

\begin{figure}[H]
\centering
\includegraphics[width=0.8\textwidth]{Gráficos/difusão_dominante.png}
\caption{Regime dominado pela difusão.}
\end{figure}

\subsection{Regime equilibrado}

Observa-se coexistência entre transporte e difusão.

\begin{figure}[H]
\centering
\includegraphics[width=0.8\textwidth]{Gráficos/equilibrado.png}
\caption{Regime equilibrado.}
\end{figure}

\subsection{Comparação final}

A análise global evidencia a influência do número de Péclet na dinâmica do sistema. Para valores elevados domina o transporte, enquanto para valores reduzidos domina a difusão.

\begin{figure}[H]
\centering
\includegraphics[width=0.8\textwidth]{Gráficos/comparacao_final.png}
\caption{Comparação final entre os diferentes regimes físicos.}
\end{figure}

\section{Conclusão}

Neste trabalho foi estudada numericamente a equação de advecção--difusão unidimensional.

Os resultados obtidos demonstram a forte influência do número de Péclet na evolução da solução, permitindo distinguir claramente regimes dominados por advecção e por difusão.

A comparação entre os esquemas numéricos evidenciou diferenças importantes em termos de precisão e difusão numérica. O esquema Upwind é mais estável, mas menos preciso, enquanto o esquema de Lax--Wendroff apresenta melhor preservação da solução.

Os resultados são consistentes com o comportamento físico esperado.

\newpage
\section*{Anexos}

\addcontentsline{toc}{section}{Anexos}

Nesta secção apresentam-se representações adicionais da solução numérica obtida para diferentes regimes físicos da equação de advecção--difusão.

\subsection*{Representações bidimensionais (2D)}

\subsubsection*{Regime dominado pela advecção}

\begin{figure}[H]
\centering
\includegraphics[width=0.85\textwidth]{Gráficos/adveccao_espaço-tempo.png}
\caption{Representação 2D da evolução no regime dominado pela advecção.}
\end{figure}

\subsubsection*{Regime dominado pela difusão}

\begin{figure}[H]
\centering
\includegraphics[width=0.85\textwidth]{Gráficos/Difusao_espaço-tempo.png}
\caption{Representação 2D da evolução no regime dominado pela difusão.}
\end{figure}

\subsubsection*{Regime equilibrado}

\begin{figure}[H]
\centering
\includegraphics[width=0.85\textwidth]{Gráficos/Equilibrado_espaço-tempo.png}
\caption{Representação 2D da evolução no regime equilibrado.}
\end{figure}

\subsection*{Representação tridimensional (3D)}

\subsubsection*{Regime dominado pela advecção}

\begin{figure}[H]
\centering
\includegraphics[width=0.85\textwidth]{Gráficos/Advecção_Mesh_3D.png}
\caption{Representação tridimensional da solução no regime dominado pela advecção.}
\end{figure}

\subsection*{Animações da simulação}

\begin{itemize}
\item \href{https://github.com/tiagobbernardo/MFE--equacao-de-adveccao-difusao/blob/7307a46e6ba824f2dffa17b3ad4cd375a18a2eb9/animations/adveccao_dominante.gif}{Regime advectivo (GIF)}
\item \href{https://github.com/tiagobbernardo/MFE--equacao-de-adveccao-difusao/blob/7307a46e6ba824f2dffa17b3ad4cd375a18a2eb9/animations/Difusao_dominante.gif}{Regime difusivo (GIF)}
\item \href{https://github.com/tiagobbernardo/MFE--equacao-de-adveccao-difusao/blob/7307a46e6ba824f2dffa17b3ad4cd375a18a2eb9/animations/equilibrado.gif}{Regime equilibrado (GIF)}
\end{itemize}

\end{document}